\documentclass[11pt, letterpaper]{article}
\usepackage[margin=1in]{geometry}
\usepackage[hyphens]{url}
\usepackage{hyperref}
\usepackage{xcolor}
\usepackage{parskip}
\usepackage{titlesec}
\usepackage{enumitem}
\usepackage{fancyhdr}

\definecolor{steelblue}{HTML}{1d4ed8}
\definecolor{mutedgray}{HTML}{6b7280}

\hypersetup{colorlinks=true, urlcolor=steelblue, linkcolor=steelblue, breaklinks=true}
\def\UrlBreaks{\do\/\do-\do_\do.}
\titleformat{\section}{\normalsize\bfseries}{}{0em}{}[\titlerule]
\pagestyle{fancy}
\fancyhf{}
\rhead{\small\color{mutedgray}Data 101 Final Project}
\lhead{\small\color{mutedgray}Earnings Volatility Analytics}
\cfoot{\small\thepage}
\renewcommand{\headrulewidth}{0.4pt}
\setlength{\parskip}{5pt}

\begin{document}

\begin{center}
  {\large\bfseries The Earnings Effect: Do Tech Stocks Actually Become More Volatile Around Earnings Time?}\\[5pt]
  {\small Data 101 · Rutgers University}\\[3pt]
  {\small Sonakshi Sharma \quad Sadhana Vasanthakumar \quad Hemadharshinii Sendhilvel \quad Rayane Skiker}\\[5pt]
  \textbf{Interactive Website:}~\href{https://sadhanavasanthakumar.github.io/data101-volatility/}{\texttt{sadhanavasanthakumar.github.io/data101-volatility}}\\
  \textbf{GitHub:}~\href{https://github.com/SadhanaVasanthakumar/data101-volatility}{\texttt{github.com/SadhanaVasanthakumar/data101-volatility}}
\end{center}

\section{The Question}

Every quarter, companies like Meta, Apple, Amazon, Netflix, and Google report
their earnings --- how much money they made, how many users they have, what
they expect going forward. These reports are called \textbf{earnings
announcements}, and the days surrounding them are informally known as
\textbf{earnings season}. Stock prices can swing wildly on these days:
Netflix lost \$50B in market value overnight in April 2022; Meta fell 26\%
in a single session in February 2022.

We asked: is this just bad luck, or do earnings announcements
\textit{reliably} make these stocks more unpredictable, every quarter, year
after year? And once a stock gets rattled by an earnings report, how long
does the turbulence actually last?

\section{The Data}

We pulled daily stock prices for META, AAPL, AMZN, NFLX, and GOOGL from
Yahoo Finance (January 2019 -- December 2024, $\approx$1,508 trading days per
stock). We manually collected 24 quarterly earnings dates per company from
Nasdaq's public calendar (120 events total) rather than using automated
databases, which are known to be unreliable for historical dates.

We measure each day's movement as its \textbf{absolute daily return} --- how
much the stock went up or down, regardless of direction. A 3\% day means the
stock moved 3\%, whether that was a gain or a loss. For each earnings date,
we flagged a 5-day window (3 days before the announcement, the day itself, and
1 day after) as ``earnings-adjacent'' --- information tends to leak before
announcements, and markets often keep reacting the day after. Everything else
was labeled ``normal trading.''

\section{What We Did --- Course Topics Used}

\textbf{Data exploration \& visualization.} We started by just looking at the
data: how much do these stocks typically move, what are their worst and best
single days, how did their prices grow relative to each other over six years.
Interactive charts for all of this are on the website.

\textbf{Hypothesis testing.} We used a \textbf{Welch's $t$-test} to formally
ask: is the average daily movement during earnings windows significantly higher
than on normal days? Result --- yes, for all five stocks, at a strict 1\%
significance level. Meta moves almost twice as much during earnings windows
(3.2\%) as on normal days (1.6\%). The test rules out that this difference is
random chance.

\textbf{Modeling.} We explored a \textbf{GARCH model} to understand how
volatility behaves over time. The core insight it captures: a volatile day
tends to be followed by another volatile day --- calm periods and stormy
periods both cluster together. By fitting this model to each stock, we found
that after an earnings shock, \textit{the turbulence doesn't just disappear}.
Netflix's elevated volatility takes roughly 95 trading days ($\approx$4.5
months) to decay to half its original size. Google recovers in about 9 days.
We also used the model to predict future volatility and tested those
predictions against what actually happened.

\textbf{Bayesian inference.} For Netflix (the most volatile stock), we ran a
Bayesian version of the same model. Instead of giving one best-guess answer,
Bayesian inference gives a \textit{range} of plausible answers and shows how
confident we should be in each. The result: forecast intervals that
actually captured 93\% of outcomes, compared to 86\% for the standard model.
Acknowledging uncertainty made the predictions more honest and more accurate.

\section{Limitations}

Our 5-day earnings window is a judgment call --- different window sizes would
shift the numbers slightly, though the core finding holds. We only studied
five large US tech stocks, so results may not generalize to other sectors or
smaller companies. The GARCH model also assumes returns behave somewhat like a
bell curve, but in reality extreme days (like Netflix $-$43\%) happen more
often than any bell-curve model predicts. These are known trade-offs in
financial modeling, not oversights.

\section{How to Explore Our Work}

\textbf{Our main deliverable is the interactive blog:}\\
\href{https://sadhanavasanthakumar.github.io/data101-volatility/}{\texttt{sadhanavasanthakumar.github.io/data101-volatility}}

This is where the full project lives --- it presents all of our findings as a
public-facing blog post with charts you can interact with: rolling volatility
over time, price growth across all five stocks, a correlation heatmap,
hypothesis test results, and the modeled risk level for each company. Each
section is written to be readable without a statistics background.

\textbf{All code and data is on GitHub:}\\
\href{https://github.com/SadhanaVasanthakumar/data101-volatility}{\texttt{github.com/SadhanaVasanthakumar/data101-volatility}}

\texttt{analysis.r} is the full R pipeline that runs every analysis;
\texttt{outputs/results.json} is the pre-computed output that powers the blog.
To reproduce from scratch: install R with packages \texttt{quantmod},
\texttt{rugarch}, \texttt{jsonlite}, \texttt{xts}, \texttt{zoo},
\texttt{moments}, then run \texttt{Rscript analysis.r}.

\end{document}